\documentclass[UTF8,11pt]{ctexart}

\usepackage[a4paper,margin=2.35cm]{geometry}
\usepackage{amsmath,amssymb,mathtools}
\usepackage{booktabs,longtable,array}
\usepackage{enumitem}
\usepackage{graphicx}
\usepackage{xcolor}
\usepackage{hyperref}

\hypersetup{
  colorlinks=true,
  linkcolor=blue!55!black,
  citecolor=blue!55!black,
  urlcolor=blue!55!black
}
\setlist{nosep,leftmargin=2em}
\newcommand{\Rhat}{\widehat{R}}
\newcommand{\Xhat}{\widehat{X}}
\newcommand{\dd}{\mathrm{d}}
\newcommand{\clade}{\mathcal{C}}

\title{仅末端量测配电网隐藏拓扑辨识\\
当前算法、实验结果与真实数据验证}
\author{terminal\_load\_only\_case33 项目技术报告}
\date{2026年7月}

\begin{document}
\maketitle

\begin{abstract}
本文整理当前项目的完整技术主线：用逐时段径向交流潮流生成量测，
从根节点参考电压与末端智能表的有功、无功和电压幅值平方估计
reduced sensitivity matrix；由对称非负的 $R/X$ 估计构造加性阻抗距离；
利用根约束 Neighbor Joining（RNJ）显式生成隐藏节点；再通过对已有噪声量测
叠加小扰动，识别稳定的根向 clade，将高置信边缘结构收缩为 pseudo node，
仅重辨识缩减后的内部主干，最后展开冻结的局部结构。

文档同时给出当前评价指标、合成算例结果、参数来源和适用边界。作为下一阶段，
项目接入公开的挪威工业配电网数据集：拓扑、支路参数和小时有功负荷来自真实且
匿名化的数据，缺失的无功和电压分别用功率因数假设及交流潮流补全。因此该结果
属于“真实拓扑与真实负荷驱动的半实证验证”，不能表述为现场 $P/Q/V$ 全量测验证。
\end{abstract}

\tableofcontents
\newpage

\section{问题定义与实现边界}

\subsection{严格末端量测场景}

设配电网为以变电站或配变低压侧为根节点 $s$ 的径向树
$T=(V,E)$。节点集合分为根节点、隐藏内部节点 $H$ 和末端用户节点 $L$：
\[
V=\{s\}\cup H\cup L.
\]
当前主场景采用如下假设：
\begin{enumerate}
  \item 根节点 $s$ 可观测，提供 $v_0(t)$ 以及必要时的总功率；
  \item 只有 $i\in L$ 具有负荷或分布式电源，并提供同步 $P_i(t),Q_i(t),V_i(t)$；
  \item 所有 $h\in H$ 为零注入、无智能表的隐藏节点；
  \item 线路闭合拓扑为单根树，支路阻抗为正；零阻抗支路需要单独讨论；
  \item 最终目标是恢复包含隐藏分支点的\emph{根向缩减树}，而不是只在末端节点
  完全图上求一棵最小生成树。
\end{enumerate}

原始 IEEE case33bw 不满足第二、三条，因为多处非叶母线带有负荷。项目的
\texttt{hybrid\_leaf} terminalization 将内部负荷迁移到新增服务线末端，
并保留原始物理叶节点，从而使内部节点严格零注入。

\subsection{当前主线与非主线方法}

需要严格区分三种对象：
\begin{itemize}
  \item \textbf{terminal-equivalent MST}：仅在观测末端之间选边，不生成隐藏节点。
  它可作距离质量诊断，但当前主线不再用它表示物理隐藏树。
  \item \textbf{RNJ/RG/NJ 隐藏树}：输入末端加性距离，输出算法生成的隐藏节点。
  当前主线采用已知根约束的 RNJ；普通 NJ 与 recursive grouping（RG）保留为基线。
  \item \textbf{候选物理图辨识}：隐藏节点和候选边由 GIS 等外部信息给定，
  只判断边状态。该任务拥有更多先验，不应与 terminal-only latent-tree 结果直接比较。
\end{itemize}

Soumalas 等假设智能表位于叶节点、根已知且中间节点不监测
\cite{soumalas2017}；Flynn 等进一步强调配变电压变化必须进入灵敏度估计
\cite{flynn2023}；Pengwah 等研究智能表覆盖不完整时的阻抗模型与两阶段定位
\cite{pengwah2024}。当前工程代码完整实现的是“根电压校正、灵敏度距离、RNJ、
扰动稳定性与一层 pseudo 收缩”这一主线。Prüfer 整数线长枚举、Flynn 原文的完整
改进 RG 与 Pengwah 的部分覆盖阻抗优化仍属于论文接口或近似复现，不能写成等价复现。

\subsection{数据生成与噪声约定}

所有当前合成算例的真值电压和线路潮流均由平衡单相等值的径向
backward--forward sweep 交流潮流逐时段计算，而非用 LinDistFlow 直接生成。
估计阶段再使用线性灵敏度模型，因此实验中保留了交流潮流与线性模型之间的系统误差。

标准噪声设置为：
\[
\sigma_P=\sigma_Q=0.5\%\times |P_i(t)|,\qquad
\sigma_V=0.02\%\times |V_i(t)|.
\]
这里电压噪声是相对于\emph{节点本地电压幅值}的标准差，并非相对于电压降或
电压差分。根节点电压维持日内公共波动，并直接作为已观测模态进入压降目标。

\section{从交流量测到阻抗距离}

\subsection{根节点参考的电压幅值平方模型}

令负荷消耗取正号。对末端节点电压幅值平方写成
\begin{equation}
  y_i(t)=v_0^2(t)-v_i^2(t).
  \label{eq:drop-target}
\end{equation}
在径向 LinDistFlow 一阶近似下，
\begin{equation}
  y(t)\approx R p(t)+X q(t)+b_s+e(t),
  \label{eq:linear-drop}
\end{equation}
其中 $s$ 表示运行场景，$b_s$ 是每个场景的固定截距，$e(t)$ 同时包含量测噪声、
交流潮流高阶项、支路损耗、相不平衡和未建模设备。

对终端 $i,j$，理想 reduced sensitivity 的元素等于根到两节点公共路径上的阻抗和：
\begin{equation}
  R_{ij}=2\sum_{e\in \operatorname{path}(s,i)\cap
  \operatorname{path}(s,j)}r_e,
  \quad
  X_{ij}=2\sum_{e\in \operatorname{path}(s,i)\cap
  \operatorname{path}(s,j)}x_e,
\end{equation}
其中因子 $2$ 对应电压平方模型；若使用电压幅值模型则不含该因子。

\subsection{多场景物理约束估计}

将多个场景按行堆叠，并为每个场景加入独立截距。代码首先解
\begin{equation}
  \min_{A,B,\{b_s\}}
  \sum_s\left\|Y_s-P_sA^\top-Q_sB^\top
  -\boldsymbol{1}b_s^\top\right\|_F^2
  +\alpha(\|A\|_F^2+\|B\|_F^2).
  \label{eq:multiscenario-regression}
\end{equation}
当 $\alpha=0$ 时直接用 SVD 最小二乘求解，不再显式形成
$X^\top X$；当 $\alpha>0$ 时用增广设计矩阵实现 Ridge。初始估计先对称化为
\begin{equation}
  \Rhat=\left[\frac{A+A^\top}{2}\right]_+,
  \qquad
  \Xhat=\left[\frac{B+B^\top}{2}\right]_+.
  \label{eq:projected-rx}
\end{equation}
对径向树还可写成
\[
R=2A_r^{-1}\operatorname{diag}(r)A_r^{-\top},
\]
故 $R/2$ 是 grounded Laplacian 的逆；只保留末端后的 $R_{LL}^{-1}$
是 Kron-reduced grounded Laplacian。由此得到
\begin{equation}
R=R^\top,\quad R\ge0,\quad R\succeq0,\quad
R_{ii}-R_{ij}\ge\epsilon,\quad R_{jj}-R_{ij}\ge\epsilon,
\label{eq:matrix-constraints}
\end{equation}
$X$ 同理。最后两个不等式正是 Pengwah 等式 (10)--(11) 使用的
shared-path 上界\cite{pengwah2024}。

代码现提供三级约束：
\begin{itemize}
  \item \texttt{basic}：对称、元素非负；
  \item \texttt{ordered}（默认）：再施加式\eqref{eq:matrix-constraints}的对角上界；
  \item \texttt{tree\_covariance}：进一步做半正定特征值投影。
\end{itemize}
初始可行投影后，代码用带回溯的 projected gradient 继续最小化原多场景目标，
每一步重新计算最优场景截距。这样比一次 Frobenius 投影更接近带约束最小二乘。
同时报告 $R^{-1}/X^{-1}$ 的 M-matrix 违例和距离四点条件残差。完整推导与消融见
\texttt{docs/sensitivity\_matrix\_constraints.md}。
固定 16 条件消融的汇总如下：
\begin{center}
\small
\begin{tabular}{lccc}
\toprule
约束模式 & basic & ordered & tree covariance \\
\midrule
平均 clade F1 & 0.8974 & 0.8630 & 0.8689 \\
平均四点违例 & 0.1396 & 0.1230 & 0.1035 \\
\bottomrule
\end{tabular}
\end{center}
因此硬约束显著改善矩阵可行性，
但在交流--线性模型误差较大时不保证提高拓扑 F1；论文结果必须同时报告两类指标。
进一步的 96 组有噪声 AC 数据条件包含 4 个算例、48/96/288 点、1/3/5/10 个场景和
2 个独立重复，共得到 288 次矩阵拟合与 2592 次 RNJ 重构。正常工况在评分前定义为
场景数不少于 3；单场景作为压力测试单独报告，不依据输出 F1 事后删样本。在固定混合
距离和容差因子 0.16 下，basic、ordered、tree covariance 的正常工况平均 clade F1
分别为 $0.9681,0.9681,0.9638$，完全恢复率分别为 $69.4\%,72.2\%,68.1\%$。
288 点且不少于 3 个场景时，三种约束平均 F1 均为 $0.9972$，完全恢复率均为
$95.8\%$；单场景压力测试不纳入上述正常指标。

同条件文献方法比较统一使用相同有噪声 AC 数据、根节点量测、预处理和 terminal-split
F1。72 个正常数据条件下，Flynn root-corrected proxy、本文 ordered RNJ、PSD RNJ
的平均 split F1 分别为 $0.9672,0.9661,0.9621$，完全恢复率分别为
$69.4\%,72.2\%,68.1\%$。Pengwah ordered-RG proxy 与 Park/Deka RG baseline
分别为 $0.6554$ 和 $0.6541$。其中 proxy 明确表示未声称复现论文缺失的整数
Prüfer 或完整迭代图学习细节；受控 NJ/RG/RNJ 行用于隔离重构算法差异。

默认时间预处理为逐日去均值：
\[
\widetilde z_s(t)=z_s(t)-\frac{1}{T_s}\sum_{\tau=1}^{T_s}z_s(\tau),
\]
以去除场景固定水平并保留日内共同激励。当前固定跨算例主线仅使用该预处理，
不再生成频域候选，也不使用真值选择滤波窗口。

\subsection{加性距离}

由 $\Rhat,\Xhat$ 构造
\begin{align}
  d^R_{ij}&=\widehat R_{ii}+\widehat R_{jj}-2\widehat R_{ij},\\
  d^X_{ij}&=\widehat X_{ii}+\widehat X_{jj}-2\widehat X_{ij}.
  \label{eq:distance}
\end{align}
在精确树模型下，$d^R_{ij}$ 等于 $i$ 到 $j$ 路径的电阻和，$d^X_{ij}$
等于电抗和。为避免 $R/X$ 绝对尺度支配组合距离，代码先以各距离矩阵正元素均值归一化：
\[
\bar d^R=\frac{d^R}{\operatorname{mean}(d^R_{ij}>0)},\qquad
\bar d^X=\frac{d^X}{\operatorname{mean}(d^X_{ij}>0)},
\]
再使用当前跨算例固定组合
\begin{equation}
  d=0.75\bar d^R+0.25\bar d^X.
  \label{eq:mixed-distance}
\end{equation}
该 $75/25$ 比例来自当前合成 case bank 的稳健性比较，不是普适物理常数。

对 RNJ 而言，无需先构造式~\eqref{eq:distance} 再把距离还原为公共路径量。
定义 $s_R,s_X$ 为 $d^R,d^X$ 的正元素均值，则当前实现直接使用
\begin{equation}
  S_{ij}=0.75\frac{\widehat R_{ij}}{s_R}
  +0.25\frac{\widehat X_{ij}}{s_X},
  \qquad H_i=S_{ii}.
  \label{eq:direct-shared-score}
\end{equation}
这里 $S_{ij}$ 就是根到 $\operatorname{LCA}(i,j)$ 的估计公共路径深度。
混合距离仍作为加性距离诊断和 NJ/RG 等方法的输入保留，但 RNJ 不再执行
``$R/X\to d\to\rho\to R/X$'' 的代数往返。

\section{已知根约束的隐藏树生成}

\subsection{直接共享路径分数}

普通 NJ 只使用叶间距离，输出无根树；事后把根挂到某个节点或某条边不能保证
恢复真实馈线根分区。对 reduced sensitivity matrix，$R_{ij}$ 与 $X_{ij}$
本身就是根到 $\operatorname{LCA}(i,j)$ 的电阻/电抗公共路径和。因此当前主线
直接使用式~\eqref{eq:direct-shared-score} 的 $S_{ij}$ 排序，并以
$H_i=S_{ii}$ 作为末端根深度。公共路径分数越大，两个 active node 越可能在
下游共享同一父节点。由同一 $S$ 构造距离后再计算
$(H_i+H_j-d_{ij})/2$ 只会原样得到 $S_{ij}$，没有额外信息，代码已删除该往返。

\subsection{RNJ 递归}

每轮在 active 集合中选取 $S_{ij}$ 最大的一对。输入时已保证
$0\le S_{ij}\le\min(H_i,H_j)$，故新隐藏父节点 $u$ 的深度直接取
$H_u=S_{ij}$，不再重复计算一个恒等的 \texttt{min}。随后添加
\[
\ell(u,i)=H_i-H_u,\qquad \ell(u,j)=H_j-H_u.
\]
若其他 active node $k$ 满足
\begin{equation}
  H_u-\min\{S_{ik},S_{jk}\}\le \tau,
\end{equation}
则将 $k$ 一并归入该父节点，允许恢复多叉树而非强制二叉化。新父节点与其余节点
的共享路径由家族成员估计的中位数更新。若最大共享路径不超过 $\tau$，剩余节点
直接连接到指定根 $s$。最后收缩非终端之间的零长度边。

容差写成
\begin{equation}
  \tau=\gamma\operatorname{median}(H_i),
\end{equation}
当前固定 $\gamma=0.16$。它来自合成 case bank 的跨算例选择，下一阶段应改为
由最小物理边长或 bootstrap 不确定度自适应确定。RNJ 的加性度量思想源自网络
层析中的根向拓扑推断\cite{ni2010,ni2011}。

\subsection{根节点保证}

RNJ 的根不是从最终无根树中“猜出”的。根标签、根到末端深度和根共享路径在每轮
递归中均参与计算，收尾时所有剩余分支只连接到给定根。因此算法保证输出树含指定根，
但不保证在灵敏度估计有偏时恢复正确的根分区。后者仍取决于 $H_i$ 与 $S_{ij}$ 的估计误差。

\subsection{可恢复性的必要边界}

若所有边权严格为正，且末端距离是精确 additive tree metric，则最小隐藏树在抑制
degree-2 隐藏节点和忽略隐藏标签置换后可唯一恢复\cite{buneman1971,semple2003}。
以下结构不能仅由末端距离唯一辨识：
\begin{itemize}
  \item degree-2 隐藏链：只能辨识链上总阻抗，无法确定中间节点个数；
  \item 零阻抗内部边：边两端诱导完全相同的距离，必须收缩；
  \item 无下游激励或严格共线注入：灵敏度列不可独立估计；
  \item 被未量测内部负荷破坏的零注入假设。
\end{itemize}

\section{稳定边收缩与一层 pseudo 主干重辨识}

\subsection{以已有噪声数据为基准的稳定性}

稳定性检验不把无噪声真值作为基准。先在实际可获得的有噪声数据 $D$ 上得到基准树
$\widehat T^{(0)}$，然后对 $D$ 额外叠加一层小扰动：
\[
\sigma_{P,Q}^{\rm add}=0.25\%,\qquad
\sigma_V^{\rm add}=0.01\%,
\]
重复估计 $B$ 次。对基准树中的 rooted clade $C$ 定义
\begin{equation}
  \pi_C=\frac{1}{B}\sum_{b=1}^{B}
  \mathbf{1}\{C\in\clade(\widehat T^{(b)})\}.
  \label{eq:bootstrap-confidence}
\end{equation}
同理可统计 sibling pair 的保持频率。当前选择规则为
$\pi_C\ge0.9$、$2\le |C|\le5$，按 clade 大小从大到小选择互不重叠的候选。
这里“稳定”表示对附加扰动不敏感，不等价于“必然正确”；系统偏差可使错误边也很稳定。

\subsection{pseudo 节点的功率与电压}

对选中的末端 clade $C$，功率严格守恒聚合：
\begin{equation}
  P_C(t)=\sum_{i\in C}P_i(t),\qquad
  Q_C(t)=\sum_{i\in C}Q_i(t).
  \label{eq:pseudo-pq}
\end{equation}
简单电压代理为
\[
\bar v_C^2(t)=\frac{1}{|C|}\sum_{i\in C}v_i^2(t).
\]
为了回推 clade 边界电压，代码先从后代 $R/X$ 行估计边界灵敏度行
$\widehat R_C,\widehat X_C$，再对每个后代构造
\begin{equation}
  \widehat v_{C,i}^2(t)=v_i^2(t)
  +\bigl(\widehat R_i-\widehat R_C\bigr)p(t)
  +\bigl(\widehat X_i-\widehat X_C\bigr)q(t).
  \label{eq:deembedding}
\end{equation}
多个后代估计取中位数得到 $\widehat v_{C,\rm deembed}^2$。由于灵敏度行相减会放大
估计误差，当前采用 shrinkage：
\begin{equation}
  \widehat v_C^2=
  \bar v_C^2+\lambda
  \left(\widehat v_{C,\rm deembed}^2-\bar v_C^2\right),
  \qquad \lambda=0.25.
  \label{eq:voltage-shrinkage}
\end{equation}

\subsection{一层收缩流程}

当前推荐流程如下：
\begin{enumerate}
  \item 用末端 $P/Q/V$ 与根电压估计 $\Rhat,\Xhat$；
  \item 构造混合阻抗距离并运行根约束 RNJ；
  \item 对已有噪声数据叠加小扰动，统计 clade/sibling 稳定度；
  \item 选择互不重叠的高置信边缘 clade，冻结其局部 rooted clades；
  \item 按式\eqref{eq:pseudo-pq}--\eqref{eq:voltage-shrinkage}构造一层 pseudo $P/Q/V$；
  \item 仅在 pseudo 节点层重新估计灵敏度和 RNJ 主干；
  \item 展开 pseudo clade，并原样放回冻结的高置信局部结构。
\end{enumerate}

局部二次重辨识有两种实验选项：以 pseudo 电压作为局部根重新回归，或直接从全局距离
中减去边界以上公共路径后 reroot。60 个独立噪声日中，冻结方案平均 F1 为
$0.9898$，distance-reroot 局部重辨识为 $0.9690$；因此局部重辨识保留为实验分支，
并未替换冻结主线。

\subsection{伪代码}

\begin{center}
\fbox{\begin{minipage}{0.93\textwidth}
\textbf{输入：} 多场景末端量测 $\{P_s,Q_s,V_s,v_{0,s}\}$，根 $s$。\\
\textbf{输出：} 根向隐藏树的 clade 集合 $\widehat{\mathcal C}$。
\begin{enumerate}
  \item 逐场景计算 $Y_s=v_{0,s}^2\mathbf1-V_s^2$ 并去均值；
  \item SVD 回归后投影得到对称非负 $\Rhat,\Xhat$；
  \item 直接由 $\Rhat,\Xhat$ 计算归一化共享路径分数 $S$ 与根深度 $H=\operatorname{diag}(S)$；
  \item 运行 RNJ 得到基准隐藏树；
  \item 对量测重复加扰并重跑步骤 1--4，估计 $\pi_C$；
  \item 收缩满足阈值的互斥边缘 clade，构造 pseudo $P/Q/V$；
  \item 在 pseudo 层运行步骤 1--4；
  \item 展开 pseudo clades，并并入冻结的局部 clades。
\end{enumerate}
\end{minipage}}
\end{center}

\section{评价指标与合成算例结论}

\subsection{rooted-clade F1}

对根向树的每个内部节点，取其下游末端集合；去除单元素集合和全体末端集合后，
得到非平凡 rooted clade 集合 $\mathcal C(T)$。隐藏节点标签不参与比较，degree-2
节点抑制后仍对应同一结构。令
\[
TP=|\widehat{\mathcal C}\cap\mathcal C^\star|,quad
FP=|\widehat{\mathcal C}\setminus\mathcal C^\star|,quad
FN=|\mathcal C^\star\setminus\widehat{\mathcal C}|.
\]
则
\begin{equation}
  \operatorname{Precision}=\frac{TP}{TP+FP},\quad
  \operatorname{Recall}=\frac{TP}{TP+FN},\quad
  F_1=\frac{2TP}{2TP+FP+FN}.
  \label{eq:clade-f1}
\end{equation}
$F_1=1$ 表示缩减根向树结构完全恢复，但不表示隐藏节点原始编号、degree-2 节点个数、
零阻抗设备层级或每条支路参数均完全恢复。

该指标与 rooted Robinson--Foulds 距离满足
\[
F_1=1-\frac{d_{RF}}{|\widehat{\mathcal C}|+|\mathcal C^\star|}.
\]
terminal-only MST 不能与 full physical edge set 直接计算边 F1；只有显式生成隐藏节点后，
才使用 rooted clade F1 或在已知候选隐藏节点图上计算物理边指标。

\subsection{当前大规模结果}

固定参数为：$P/Q$ 噪声 $0.5\%$、电压噪声 $0.02\%$、逐日去均值、
$0.75d_R+0.25d_X$、RNJ 容差因子 $0.16$、稳定阈值 $0.9$、
最大收缩 clade 大小 $5$、pseudo 电压去嵌入权重 $0.25$。

\begin{table}[htbp]
\centering
\caption{一层冻结聚合在合成 case bank 上的完整 rooted-clade 指标}
\begin{tabular}{lcccc}
\toprule
实验集合 & 条件数 & 基线平均 F1 & 聚合平均 F1 & 完全恢复率变化\\
\midrule
不同场景数与数据量 & 48 & 0.8847 & 0.9073 & 56.3\% $\to$ 72.9\%\\
独立 288 点噪声日 & 60 & 0.9510 & 0.9898 & 58.3\% $\to$ 88.3\%\\
\bottomrule
\end{tabular}
\end{table}

48 个数据量条件中，聚合改善 12 个、持平 35 个、变差 1 个；60 个独立日中改善
20 个、持平 40 个、无变差。结果支持“冻结高置信边缘 clade、只重辨识缩减主干”
这一设计，但不能据此宣称所有新馈线都接近完全恢复。参数仍来自当前 case bank，
面对分布外数据时需要无标签自适应规则。

\subsection{主要误差来源}

\begin{enumerate}
  \item \textbf{激励与条件数：} 住宅负荷共同日周期、工业负荷共线以及固定功率因数
  使 $P/Q$ 设计矩阵接近低秩。
  \item \textbf{量测噪声相对差分放大：} $V$ 的相对误差虽小，但差分、去公共模态和
  灵敏度回归使用的是更小的电压变化量。
  \item \textbf{交流--线性模型偏差：} 损耗和二阶电压项使式\eqref{eq:linear-drop}
  的误差具有结构性，并非独立白噪声。
  \item \textbf{距离竞争 margin 小：} 多个内部 split 的候选加性距离接近时，
  小偏差会改变 RNJ 最大共享路径对。
  \item \textbf{不可辨识结构：} degree-2 链和零阻抗边没有末端距离签名。
  \item \textbf{稳定但错误：} bootstrap 只检验局部扰动稳定性，不能消除固定系统偏差。
\end{enumerate}
\subsection{RNJ 前后的错误区域变化}

为检验早期“错误主要位于中心内部节点”的结论，NJ、RG、RNJ 在 72 个正常条件下使用
完全相同的 ordered $R/X$ 估计和 $0.75d_R+0.25d_X$ 距离。结果如下：
\begin{center}
\begin{tabular}{lrrrr}
\toprule
方法 & 平均 split F1 & 完全恢复率 & 漏真实 split & 多预测 split\\
\midrule
NJ  & 0.6298 & 0.0\%  & 1  & 469\\
RG  & 0.7417 & 4.2\%  & 15 & 268\\
RNJ & 0.9661 & 72.2\% & 3  & 26\\
\bottomrule
\end{tabular}
\end{center}
RNJ 对根邻接中心主干的 198 次识别全部正确；RG 漏 4 次，NJ 虽未漏真实中心 split，
但生成大量额外二叉 split。RNJ 剩余 26 条额外 split 中，20 条只包含不超过 3 个
terminal，表现为末端簇内部过分裂。因此 RNJ 并非以牺牲边缘为代价改善中心，而是同时
压低两类误差；中心误差消失后，残余误差的组成自然偏向边缘。早期“内部错误”标签把
所有额外隐藏 split 都记为内部边，即使它仅切分一个很小的末端簇，因此物理位置含义过粗。

\section{公开真实数据调研与首次辨识}

\subsection{数据源筛选}

\begin{longtable}{p{3.1cm}p{3.2cm}p{4.1cm}p{4.2cm}}
\caption{与末端量测拓扑辨识相关的开放数据源}\label{tab:datasets}\\
\toprule
数据源 & 规模 & 可用信息 & 与当前方法的关系\\
\midrule
\endfirsthead
\toprule
数据源 & 规模 & 可用信息 & 与当前方法的关系\\
\midrule
\endhead
挪威工业 MV/LV，v2.3\cite{sandell2023}
& 两条径向，45 个客户，超过三年小时数据
& 匿名化真实拓扑、公共标幺支路参数、客户有功电量
& 已接入；无现场 $Q/V$，只能做真实负荷驱动的半实证验证\\
Iowa Distribution Test Systems\cite{iowa}
& 3 条馈线、240 节点、1120 客户、2017 全年
& 真实拓扑和元件参数；二次变压器层聚合小时电量
& 适合下一批中等规模实验；同样缺少逐客户同步 $Q/V$\\
Non-Synthetic European LV\cite{arboleya2019}
& 10290 节点、8087 客户、20 天
& 真实 GIS、智能表读数、OpenDSS 构建与潮流脚本
& 规模大且三相细节丰富；需先解析数据许可与相别/中性线模型\\
挪威城乡 LV reference dataset\cite{norwaylv2024}
& 约 1.2 GB，多组城乡低压网
& 网架 CSV 与负荷数据
& 适合扩展北欧低压场景，尚未接入\\
SMART-DS\cite{smartds}
& 可达数千馈线、百万客户
& 完整合成网、负荷、DER、天气与多场景
& 不是实测数据；适合扩展性和分布外压力测试\\
SimBench\cite{simbench}
& LV 至 EHV 的标准网与全年时序
& 网络、负荷、发电和储能曲线
& 真实感基准而非现场 PQV；适合可重复跨电压等级测试\\
\bottomrule
\end{longtable}

公开数据的核心缺口非常明确：同时公开“低压客户级拓扑真值、同步 $P/Q/V$、根电压、
足够长时间序列”的数据极少。许多真实 AMI 数据只提供计费用有功电量，而真实电压数据
常因隐私和资产安全无法与完整拓扑共同发布。

\subsection{挪威数据的适配}

项目采用 Zenodo v2.3 数据\cite{sandell2023}。数据由 Norgesnett 与 CINELDI 项目
采集后匿名化、简化并公开，许可证为 CC BY 4.0。适配器
\texttt{terminal\_case33/data/norwegian\_industrial.py} 完成：
\begin{enumerate}
  \item 按 47~kV 根节点分离两条闭合径向分量；
  \item 剪除没有智能表负荷的死端，但不删除任何通向测量客户的路径；
  \item 将字符串母线名映射为整数，保留反向映射和来源元数据；
  \item 将发布的公共标幺 $R/X$ 保持不变，并在 $S_{base}=1$~MVA、
  $V_{base}=11$~kV 等值基准上转换为欧姆，以复用现有交流潮流器；
  \item 保留零阻抗理想变压器/连接支路，明确标记为公共标幺平衡等值模型。
\end{enumerate}

两条径向在剪枝后分别为 54 节点/40 个末端表和 17 节点/5 个末端表，真实客户负荷
本来就位于叶节点，无需像 case33 那样迁移负荷。

\subsection{半实证量测构造}

源文件只提供小时有功电量。每小时的 \texttt{Load\_kWh} 数值按该小时平均 kW 使用。
为运行 AC 潮流，统一假设功率因数 $0.95$：
\[
Q_i(t)=P_i(t)\tan(\arccos 0.95).
\]
根电压设为
\[
v_0(t)=1.02+0.0015\sin(2\pi t/24).
\]
逐时段 AC 潮流生成末端真值电压，再加入 $P$ 相对噪声 $0.5\%$ 和相对本地电压幅值
$0.02\%$ 的电压噪声。由于 $Q$ 不是实测，辨识器采用 active-only 模式：
回归中 $Q$ 列置零，$\Rhat$ 实际吸收了固定功率因数下的有效
$R+\tan(\phi)X$ 距离。该距离仍为加性距离，但不能解释为独立恢复的物理 $R$。

\subsection{零阻抗 clade 与可辨识真值}

发布模型含大量零阻抗理想连接。若内部边 $e=(u,v)$ 满足 $r_e+x_e=0$，
收缩前后所有末端距离完全相同。因此实验同时定义：
\begin{itemize}
  \item \textbf{physical clade}：发布树中所有内部节点的下游末端集合；
  \item \textbf{impedance-identifiable clade}：只有父边阻抗严格为正的非平凡下游集合。
\end{itemize}
径向 1 有 6 个 physical clades，但仅 3 个阻抗可辨识 clades；径向 2 分别为 2 和 1。
即使距离完全准确，若直接和未收缩 physical clades 比较，F1 上限也只有 $2/3$。

\subsection{首次结果}

\begin{table}[htbp]
\centering
\caption{真实拓扑与真实有功负荷驱动的 active-only RNJ 结果}
\begin{tabular}{cccccccc}
\toprule
径向 & 末端数 & 天数 & 样本数 & 可辨识 clade & 基线 F1 & 聚合 F1 & 完全恢复\\
\midrule
1 & 40 & 7  & 168  & 3 & 0.000 & 0.000 & 否\\
1 & 40 & 30 & 720  & 3 & 0.000 & 0.000 & 否\\
1 & 40 & 90 & 2160 & 3 & 0.000 & 0.000 & 否\\
2 & 5  & 7  & 168  & 1 & 1.000 & 1.000 & 是\\
2 & 5  & 30 & 720  & 1 & 1.000 & 1.000 & 是\\
2 & 5  & 90 & 2160 & 1 & 1.000 & 1.000 & 是\\
\bottomrule
\end{tabular}
\end{table}

径向 2 的唯一正阻抗 split 在所有窗口均完全恢复。径向 1 则完全失败，并且未选出
满足 $0.9$ 稳定阈值的 pseudo cluster。该失败不是“样本越多反而必然更差”，而是
以下因素共同作用：
\begin{enumerate}
  \item 40 个客户只对应 3 个正阻抗可辨识内部 split，网络含大规模零阻抗多叉层；
  \item 多个同一 11~kV 接入点的客户在模型中电压真值相同，独立电压噪声却直接进入
  40 维回归，逐客户 service sensitivity 实际为零；
  \item 工业负荷幅值跨越多个数量级且具有共同生产周期，active-only 设计矩阵弱激励；
  \item 严格 clade F1 只接受整个下游集合完全相等，接近的 34 节点大 clade
  只要多一个或少一个末端就记为一个 $FP$ 加一个 $FN$。
\end{enumerate}

因此该实验给出的结论不是“算法已经通过真实系统验证”，而是：
小规模正阻抗 split 能恢复；大规模、零 service 阻抗、仅有功且相关的工业 AMI
数据明显超出当前逐客户灵敏度 RNJ 的稳定范围。

\begin{figure}[htbp]
\centering
\includegraphics[width=0.95\textwidth]{../../outputs/norwegian_real_load_identification/radial_2_topology_comparison.png}
\caption{挪威径向 2：发布物理树与 90 天 active-only RNJ 距离树。算法树已收缩
零长度隐藏边，因此节点数不应与发布设备层逐一比较。}
\end{figure}

\section{下一阶段技术路线}

\subsection{针对真实 AMI 的优先改进}

\begin{enumerate}
  \item \textbf{先做零阻抗/同电压接入点辨识。}
  对去根公共模态后的电压序列计算
  \[
  \delta_{ij}^{V}=
  \sqrt{\frac{1}{T}\sum_t
  \left[(v_i^2-v_0^2)-(v_j^2-v_0^2)\right]^2}.
  \]
  当 $\delta_{ij}^{V}$ 与量测噪声基线一致时，先把客户聚合为 voltage-equivalent
  injection group，再对组级总 $P/Q$ 估计主干。这比要求零 service impedance
  的每个客户都独立估计灵敏度更符合挪威数据结构。

  \item \textbf{用无标签阈值替代真值调参。}
  RNJ 容差应由距离 bootstrap 标准差、物理最小正阻抗下界或稳定平台选择；
  pseudo 阈值可用 stability selection 的期望错误发现数控制。

  \item \textbf{处理 errors-in-variables。}
  $P/Q$ 也含噪声时普通最小二乘有衰减偏差。应比较 total least squares、
  instrumental variables、SIMEX 和基于量测方差的广义最小二乘。

  \item \textbf{先投影到加性树度量。}
  在 RNJ 前求解最近树度量或四点条件正则：
  \[
  \min_{D\in\mathcal T}
  \|W\odot(D-\widehat D)\|_F^2+
  \lambda\Omega(D),
  \]
  其中 $\mathcal T$ 是加性树度量集合，$W$ 来自 bootstrap 不确定度。

  \item \textbf{用 AC 残差二次筛选。}
  RNJ 生成若干稳定候选 split 后，回到 AC 潮流或 branch-flow 模型，
  用跨场景电压重构误差、线路损耗和参数正值约束重新排序，而不是只依赖线性距离。
\end{enumerate}

\subsection{真实数据扩展顺序}

建议按以下顺序推进：
\begin{enumerate}
  \item 完成挪威径向 1 的 voltage-equivalent grouping，并在不看拓扑标签的条件下
  固定阈值，再锁定一次测试集；
  \item 接入 Iowa 的 3 条馈线，复用真实拓扑和变压器层聚合负荷，验证中等规模；
  \item 从 Non-Synthetic European LV 数据中抽取单配变子树，保留三相和中性线，
  比较正序等值与三相模型；
  \item 用 SimBench 和 SMART-DS 做规模、DER 渗透率和量测缺失压力测试；
  \item 最终寻找或与 DSO 合作获得同步客户 $P/Q/V$ 与配变电压，完成真正现场验证。
\end{enumerate}

\subsection{复现实验命令}

\begin{verbatim}
# 下载约 44.5 MB 的公开数据
python scripts/download_norwegian_industrial.py

# 两条径向、7/30/90 天、带进度可复现实验
python -m experiments.run_norwegian_real_load_identification

# 合成 case bank 的完整 rooted-clade F1
python -m experiments.run_complete_rooted_topology_f1

# 编译本文档
cd docs/current_algorithm_summary
xelatex main.tex
bibtex main
xelatex main.tex
xelatex main.tex
\end{verbatim}

真实数据实验输出目录为
\nolinkurl{outputs/norwegian_real_load_identification/}。主要文件包括
\nolinkurl{identification_results.csv}、\nolinkurl{clade_details.csv}、
\nolinkurl{data_summary.csv}、\nolinkurl{metrics.json}、报告和两条径向的拓扑对比图。

\subsection{当前结论}

当前算法已经形成可运行、可评价、根约束明确的初步主线。在合成 terminal-only
case bank 上，一层高置信边缘收缩显著提高了平均 F1 和完全恢复率；但真实工业
AMI 试验表明，零阻抗设备层、仅有功量测和强相关负荷会使逐客户灵敏度距离失效。
下一步不应继续只在原有合成算例上微调滤波窗口，而应围绕“可辨识真值定义、
同电压客户预聚合、errors-in-variables 与 AC 候选重评分”推进。


\bibliographystyle{IEEEtran}
\bibliography{references}

\end{document}
